% Image credits for the photographs and AI illustrations of Book 1 (Primary and Middle School).
% Machine-readable license records live in images/book1/CREDITS.md; keep the
% two files in sync. The Book 1 writer owns this file: add one \item per
% photograph (inside the itemize, by chapter) as photographs land.
% Spanish edition: every author, licence and source of the English list is
% kept exactly; only the connecting words are translated.
\chapter*{Créditos de las imágenes}

Los dibujos, mapas y gráficos de este libro son originales; los mapas y
gráficos trazados a partir de datos publicados citan su fuente en el pie de
figura o en esta lista. Las fotografías se utilizan con sus respectivas
licencias libres, con nuestro agradecimiento a sus autores; los enlaces
completos a las fuentes y las URL de las licencias de cada fotografía
figuran en \texttt{images/book1/CREDITS.md}, en el repositorio del libro.

\bigskip
Junto a las fotografías y a los dibujos originales, algunas figuras son
ilustraciones generadas por IA, producidas con la generación de imágenes de
OpenAI a partir de instrucciones escritas por los editores del libro, y
revisadas en cuanto a su exactitud científica antes de incluirlas. Solo
muestran escenas genéricas, nunca un lugar real, una persona real ni un
ejemplar real. La instrucción completa de cada imagen generada está
registrada en \texttt{images/book1/ai/PROMPTS.md}, en el repositorio.

\bigskip
\noindent Fotografías y datos, por capítulo:
\begin{itemize}\small
  \item Cap.~\ref{ch:g2:water-on-earth} (el agua en la Tierra): la Tierra vista desde el Apolo~17, el 7~de diciembre de 1972 (``The Blue Marble''); NASA, dominio público (vía Wikimedia Commons). Los 100 cuadrados se calculan a partir del relieve NOAA ETOPO1; los 100 cubos, a partir de la tabla de la U.S. Geological Survey Water Science School (según Shiklomanov, 1993).
  \item Cap.~\ref{ch:g3:stones-pebbles-sand} (piedras, cantos rodados y arena): un muro de piedra seca de granito; fotografía de Tony Atkin, Geograph, CC BY-SA 2.0 (vía Wikimedia Commons). Tejados de losas de piedra en Suiza; fotografía de Michael C. Rygel, CC BY-SA 3.0 (vía Wikimedia Commons).
  \item Cap.~\ref{ch:g3:continents-and-oceans} (el globo terráqueo): los dos globos y el mapamundi se dibujan a partir de las costas de Natural Earth (dominio público).
  \item Cap.~\ref{ch:g4:clouds-rain-snow} (nubes, lluvia y nieve): cúmulos sobre Kansas (S.~Corfidi), manchas de cirros (R.~F.~Kresge), una capa de estratos alrededor de unas cumbres, una tormenta sobre el desierto de Arizona (colección de B.~Hooke), un granizo muy grande; todas de la NOAA Photo Library, dominio público (vía Wikimedia Commons).
  \item Cap.~\ref{ch:g4:volcanoes-earthquakes-first-look} (volcanes y terremotos): lava que entra en el océano en el Kilauea, USGS Hawaiian Volcano Observatory; el monte Santa Helena el 18~de mayo de 1980, fotografía de Austin Post, USGS; una casa arrancada de sus cimientos por el terremoto de Loma Prieta de 1989, fotografía de J.~K.~Nakata, USGS; todas de dominio público (vía Wikimedia Commons). El mapa usa la lista de volcanes holocenos del Global Volcanism Program de la Smithsonian y el catálogo de terremotos de la U.S. Geological Survey.
  \item Cap.~\ref{ch:g4:world-physical-map} (el mapa físico): el Himalaya y la Amazonia fotografiados desde la Estación Espacial Internacional, NASA, dominio público (vía Wikimedia Commons). El mapa físico se dibuja a partir del relieve NOAA ETOPO1, con ríos, lagos y costas de Natural Earth (dominio público).
  \item Cap.~\ref{ch:g5:water-shapes-the-land} (el agua modela el terreno): el río Mamoré fotografiado desde la Estación Espacial Internacional (ISS007-E-10797) y el delta del Nilo visto por el instrumento MODIS del satélite Terra, NASA, dominio público (vía Wikimedia Commons).
  \item Cap.~\ref{ch:g5:weather-and-climate} (el tiempo y el clima): el desierto de Sonora, U.S. National Park Service, Saguaro National Park, dominio público (vía Wikimedia Commons). Los climogramas se dibujan a partir de las NOAA NCEI U.S. Climate Normals 1991--2020 (dominio público).
  \item Cap.~\ref{ch:g5:natural-hazards} (los peligros naturales): el ojo del huracán Isabel fotografiado desde la Estación Espacial Internacional por Ed Lu, NASA, dominio público; el deslizamiento de tierra de 2014 en el estado de Washington, Mark Reid, U.S. Geological Survey, dominio público (ambas vía Wikimedia Commons).
  \item Cap.~\ref{ch:g6:earth-a-planet} (la Tierra, un planeta): Venus desde la Mariner~10 (NASA/JPL-Caltech); la Tierra desde el Apolo~17 (NASA); \emph{Earthrise} desde el Apolo~8 (NASA, Bill Anders); Marte y Júpiter desde el telescopio espacial Hubble (NASA, ESA, the Hubble Heritage Team y A.~Simon); todas de dominio público (vía Wikimedia Commons). El \emph{Pale Blue Dot} visto por la Voyager~1 en 1990, versión de 2020 (NASA/JPL-Caltech, PIA23645, dominio público). Datos de los planetas según la NASA planetary fact sheet.
  \item Cap.~\ref{ch:g6:latitude-longitude} (latitud y longitud): el H4 de Harrison en el National Maritime Museum, Greenwich, fotografía de Colonel Warden, CC BY-SA 3.0 (vía Wikimedia Commons). Costas y ciudades de Natural Earth (dominio público).
  \item Cap.~\ref{ch:g6:rocks-and-minerals} (rocas y minerales): granito, basalto, caliza fosilífera, arenisca, mármol, lutita, un cristal de cuarzo y una lámina de mica moscovita, fotografías de James St.~John, CC BY 2.0; un tejado de pizarra, fotografía de Øyvind Holmstad, CC0 (todas vía Wikimedia Commons).
  \item Cap.~\ref{ch:g6:inside-the-earth} (el interior de la Tierra): Inge Lehmann en 1932, fotografía de Even Neuhaus, Royal Library of Denmark, dominio público (vía Wikimedia Commons). Las capas, densidades y velocidades de las ondas se dibujan a partir del Preliminary Reference Earth Model (A.~M. Dziewonski y D.~L. Anderson, 1981).
  \item Cap.~\ref{ch:g6:water-underground} (el agua bajo tierra): Big Spring, Misuri, fotografía de Johnston9494, CC BY-SA 4.0 (vía Wikimedia Commons). Porosidades según R.~C. Heath, \emph{Basic Ground-Water Hydrology}, U.S. Geological Survey (1983); proporciones de agua dulce de la USGS Water Science School.
  \item Cap.~\ref{ch:g7:plate-tectonics} (la tectónica de placas): los continentes del sur reunidos, con las franjas de fósiles de \emph{Cynognathus}, \emph{Lystrosaurus}, \emph{Mesosaurus} y \emph{Glossopteris}; U.S. Geological Survey, \emph{This Dynamic Earth}, dominio público.
  \item Cap.~\ref{ch:g7:plate-tectonics}: la falla de San Andrés a través de la Carrizo Plain; fotografía de Robert E. Wallace, U.S. Geological Survey, dominio público.
  \item Cap.~\ref{ch:g7:plate-tectonics}: límites de placas según P.~Bird (2003), GeoJSON de H.~Ahlenius (ODC-By 1.0); costas de Natural Earth (dominio público); terremotos del catálogo de la U.S. Geological Survey; volcanes del Global Volcanism Program de la Smithsonian.
  \item Cap.~\ref{ch:g7:volcanoes} (los volcanes): el monte Santa Helena desde Johnston Ridge el 17~de mayo y el 10~de septiembre de 1980, fotografías de Harry Glicken, U.S. Geological Survey; una fuente de lava en el Kilauea, 1983, y el Mauna Loa desde el aire, 1985, fotografías de J.~D. Griggs, U.S. Geological Survey; la columna eruptiva del monte Pinatubo, 12~de junio de 1991, fotografía de Dave Harlow, U.S. Geological Survey; todas de dominio público (vía Wikimedia Commons). El mapa y los recuentos usan la lista de volcanes holocenos del Global Volcanism Program de la Smithsonian; límites de placas según P.~Bird (2003).
  \item Cap.~\ref{ch:g7:earthquakes} (los terremotos): un pueblo cerca de la costa de Sumatra después del tsunami del 26~de diciembre de 2004, fotografía de la U.S. Navy de Philip A. McDaniel; los escarpes de falla del terremoto de Borah Peak de 1983, fotografía de R.~E. Wallace, U.S. Geological Survey; la carretera hundida en el lago Hebgen, 1959, fotografía de I.~J. Witkind, U.S. Geological Survey; todas de dominio público (vía Wikimedia Commons). Los recuentos y las profundidades usan el catálogo de terremotos de la U.S. Geological Survey; la tabla de intensidades es una versión abreviada de la escala de Mercalli modificada de la U.S. Geological Survey.
  \item Cap.~\ref{ch:g7:sediments-and-fossils} (sedimentos y fósiles): el Gran Cañón desde Hermits Rest, fotografía del NPS de Kristen M. Caldon, dominio público; un ammonites, un trilobites y unas marcas de oleaje en el Capitol Reef National Park, fotografías de James St.~John, CC BY 2.0; huellas de dinosaurio en Dinosaur Ridge, fotografía de Dexter Perkins, CC0 (todas vía Wikimedia Commons).
  \item Cap.~\ref{ch:g7:geological-time} (las edades de la Tierra): James Hutton, retrato de Henry Raeburn, dominio público (vía Wikimedia Commons). Las edades y los colores de la escala de tiempo son los de la Tabla Cronoestratigráfica Internacional (International Commission on Stratigraphy).
  \item Cap.~\ref{ch:g8:the-atmosphere} (la atmósfera): lanzamiento de un globo meteorológico en Topeka, Kansas, fotografía de Stephen Corfidi, NOAA/National Weather Service; el borde de la Tierra al atardecer desde la Estación Espacial Internacional, NASA, ISS Expedition~23; ambas de dominio público (vía Wikimedia Commons). Los perfiles se calculan a partir de la \emph{U.S. Standard Atmosphere, 1976} (NOAA, NASA y U.S. Air Force).
  \item Cap.~\ref{ch:g8:winds-global-circulation} (los vientos y la circulación general): las Américas vistas por el satélite GOES-16 el 26~de julio de 2022, NOAA, dominio público (vía Wikimedia Commons). Las costas del mapa de los alisios son de Natural Earth (dominio público).
  \item Cap.~\ref{ch:g8:fronts-and-weather-maps} (masas de aire, frentes y mapas del tiempo): una borrasca frente a Islandia vista por el satélite Terra, NASA/GSFC, MODIS Rapid Response Team, Jacques Descloitres, dominio público (vía Wikimedia Commons). Costas del mapa de las masas de aire de Natural Earth (dominio público).
  \item Cap.~\ref{ch:g8:the-oceans} (los océanos): el HMS \emph{Challenger} preparándose para un sondeo, 1872, grabado de los informes de la expedición del Challenger reproducido en W.~A. Herdman, \emph{Founders of Oceanography} (1923); la corriente del Golfo vista por el instrumento MODIS del satélite Aqua, NASA Earth Observatory; ambos de dominio público (vía Wikimedia Commons). Los perfiles del fondo oceánico se dibujan a partir del relieve NOAA ETOPO1; costas de los mapas de los giros y de las boyas de Natural Earth (dominio público).
  \item Cap.~\ref{ch:g8:ice-on-earth} (el hielo en la Tierra): el glaciar Pedersen en 1990 y en 2016, fotografías de M.~Tetreau y D.~Kurtz, U.S. National Park Service; el borde de una plataforma de hielo frente a la península Antártica, NASA, Operation IceBridge (M.-J. Viñas); todas de dominio público (vía Wikimedia Commons). El gráfico del hielo marino usa el NSIDC Sea Ice Index, versión~4 (F.~Fetterer y otros, 2025, doi:10.7265/a98x-0f50).
  \item Cap.~\ref{ch:g9:climate-past-and-present} (el clima de la Tierra, ayer y hoy): un tramo del testigo de hielo GISP2 de Groenlandia, NOAA Paleoclimatology, dominio público (vía Wikimedia Commons). Los registros se dibujan a partir de NASA GISTEMP versión~4, del registro de dióxido de carbono de Mauna Loa (NOAA Global Monitoring Laboratory y Scripps Institution of Oceanography), de la temperatura de EPICA Domo~C (J.~Jouzel y otros, 2007), de la serie compuesta de dióxido de carbono de los testigos de hielo antárticos (B.~Bereiter y otros, 2015) y del nivel del mar del NOAA Laboratory for Satellite Altimetry; las afirmaciones de la evaluación son las del Sexto Informe de Evaluación del IPCC, Grupo de Trabajo~I, Resumen para responsables de políticas (2021).
  \item Cap.~\ref{ch:g9:natural-risks} (vivir con los riesgos naturales): el mapa de peligrosidad del Nevado del Ruiz, según T.~L. Wright y T.~C. Pierson, \emph{Living with Volcanoes}, U.S. Geological Survey Circular 1073 (1992); una boya de alerta de tsunamis inspeccionada por la tripulación del USCGC \emph{William Hart}, U.S. Coast Guard, 2021; ambas de dominio público (vía Wikimedia Commons).
  \item Cap.~\ref{ch:g9:earth-resources} (los recursos de la Tierra): la mina de Chuquicamata vista por ASTER, NASA/METI/AIST/Japan Space Systems y U.S./Japan ASTER Science Team, 2007, dominio público; una capa de carbón en un talud de carretera del condado de Belmont, Ohio, fotografía de James St.~John, CC BY 2.0 (ambas vía Wikimedia Commons). Los factores de concentración son de D.~Perkins y otros, \emph{Mineralogy}, y de S.~Earle, \emph{Environmental Geology}.
  \item Cap.~\ref{ch:g9:landforms} (las formas del relieve): las Seven Sisters, fotografía de Arild Vågen, 2017, CC BY-SA 4.0; dunas en media luna en la costa del sur de Brasil, fotografía de un astronauta de la NASA desde la Estación Espacial Internacional, 2017, dominio público; estalactitas en la cueva de Carlsbad (Carlsbad Cavern), fotografía del National Park Service de Daniel Leifheit, 2018, dominio público (todas vía Wikimedia Commons).
  \item Cap.~\ref{ch:g9:maps-of-the-earth} (mapas de toda la Tierra): el delta del río Lena, imagen Landsat de 2000, USGS EROS Data Center vía NASA Visible Earth, dominio público; un satélite GPS, recreación artística de la NASA, dominio público (ambas vía Wikimedia Commons). Costas de los mapamundis de Natural Earth (dominio público).
\end{itemize}
